% Part: first-order-logic
% Chapter: syntax-and-semantics
% Section: first-order-languages

\documentclass[../../../include/open-logic-section]{subfiles}

\begin{document}

\olfileid{fol}{syn}{fol}

\olsection{Eerste-ordetalen}


Uitdrukkingen van de eerste-ordelogica worden opgebouwd uit een
basisvocabulaire met \emph{!!{variable}s}, \emph{!!{constant}s},
\emph{!!{predicate}s} en soms \emph{!!{function}s}. Daaruit worden,
samen met logische connectieven, kwantoren en leestekens zoals haakjes
en komma's, \emph{termen} en \emph{!!{formula}s} gevormd.

\begin{explain}
Informeel zijn !!{predicate}s namen voor eigenschappen en relaties,
!!{constant}s namen voor afzonderlijke objecten en !!{function}s namen
voor functies. Deze tekens vormen, met uitzondering van de
!!{identity}~$\eq$, de \emph{niet-logische symbolen} en bepalen samen
een taal. Iedere eerste-ordetaal~$\Lang L$ wordt door haar
niet-logische symbolen bepaald. In het algemeenste geval bevat
$\Lang L$ van elke soort oneindig veel symbolen.
\end{explain}

In het algemene geval gebruiken we in de eerste-ordelogica de volgende
symbolen:

\begin{enumerate}
\item Logische symbolen
\begin{enumerate}
\item Logische connectieven:
  \startycommalist
  \iftag{prvNot}{\ycomma $\lnot$ (negatie)}{}%
  \iftag{prvAnd}{\ycomma $\land$ (conjunctie)}{}%
  \iftag{prvOr}{\ycomma $\lor$ (disjunctie)}{}%
  \iftag{prvIf}{\ycomma $\lif$ (!!{conditional})}{}%
  \iftag{prvIff}{\ycomma $\liff$ (!!{biconditional})}{}%
  \iftag{prvAll}{\ycomma $\lforall$ (universele kwantor)}{}%
  \iftag{prvEx}{\ycomma $\lexists$ (existentiële kwantor)}{}.
\tagitem{prvFalse}{De propositionele constante voor !!{falsity}~$\lfalse$.}{}
\tagitem{prvTrue}{De propositionele constante voor !!{truth}~$\ltrue$.}{}
\item De tweeplaatsige !!{identity}~$\eq$.
\item Een !!{denumerable}s verzameling !!{variable}s: $\Obj v_0$, $\Obj v_1$, $\Obj
  v_2$, \dots
\end{enumerate}
\item Niet-logische symbolen, die de \emph{standaardtaal} van de
  eerste-ordelogica vormen
\begin{enumerate}
\item Een !!{denumerable}s verzameling $n$-plaatsige !!{predicate}s voor elke $n>0$: $\Obj
  A^n_0$, $\Obj A^n_1$, $\Obj A^n_2$, \dots
\item Een !!{denumerable}s verzameling !!{constant}s: $\Obj c_0$, $\Obj c_1$, $\Obj
  c_2$, \dots.
\item Een !!{denumerable}s verzameling $n$-plaatsige !!{function}s voor elke $n>0$:
  $\Obj f^n_0$, $\Obj f^n_1$, $\Obj f^n_2$, \dots
\end{enumerate}
\item Leestekens: (, ) en de komma.
\end{enumerate}

De meeste definities en resultaten formuleren we voor de volledige
standaardtaal van de eerste-ordelogica. Afhankelijk van de toepassing
kunnen we de taal echter ook beperken tot slechts enkele
!!{predicate}s, !!{constant}s en !!{function}s.

\begin{ex}
De taal~$\Lang L_A$ van de rekenkunde bevat één tweeplaatsig
!!{predicate}~$<$, één !!{constant}~$\Obj 0$, één eenplaatsige
!!{function}~$\prime$ en twee tweeplaatsige !!{function}s~$+$ en~$\times$.
\end{ex}

\begin{ex}
De taal van de verzamelingenleer~$\Lang L_Z$ bevat uitsluitend het ene
tweeplaatsige !!{predicate}~$\in$.
\end{ex}

\begin{ex}
De taal van ordeningen~$\Lang L_\le$ bevat uitsluitend het
tweeplaatsige !!{predicate}~$\le$.
\end{ex}

Ook dit zijn conventies: formeel zijn het slechts alternatieve namen.
Zo zijn $<$, $\in$ en $\le$ alternatieve namen voor $\Obj A^2_0$, staat
$\Obj 0$ voor $\Obj c_0$, $\prime$ voor $\Obj f^1_0$, $+$ voor
$\Obj f^2_0$ en $\times$ voor $\Obj f^2_1$.

\iftag{defNot,defOr,defAnd,defIf,defIff,defTrue,defFalse,defEx,defAll}{%
Naast de hierboven ingevoerde primitieve connectieven en
\iftag{notprvEx,notprvAll}{kwantor}{kwantoren} gebruiken we ook de
volgende \emph{gedefinieerde} symbolen:
\startycommalist
  \iftag{defNot}{\ycomma $\lnot$ (negatie)}{}%
  \iftag{defAnd}{\ycomma $\land$ (conjunctie)}{}%
  \iftag{defOr}{\ycomma $\lor$ (disjunctie)}{}%
  \iftag{defIf}{\ycomma $\lif$ (!!{conditional})}{}%
  \iftag{defIff}{\ycomma $\liff$ (!!{biconditional})}{}%
  \iftag{defAll}{\ycomma $\lforall$ (universele kwantor)}{}%
  \iftag{defEx}{\ycomma $\lexists$ (existentiële kwantor)}{}%
  \iftag{defFalse}{\ycomma !!{falsity}~$\lfalse$}{}%
  \iftag{defTrue}{\ycomma !!{truth}~$\ltrue$}}{}.

\begin{tagblock}{defNot,defOr,defAnd,defIf,defIff,defTrue,defFalse,defEx,defAll}
\begin{explain}
Een gedefinieerd symbool behoort formeel niet tot de taal, maar wordt
als informele afkorting ingevoerd. Daarmee kunnen we formules afkorten
die met uitsluitend primitieve symbolen tamelijk lang zouden worden.
Dat is een duidelijk voordeel. Het grotere voordeel is echter dat
bewijzen korter worden. Ieder primitief symbool moet in bewijzen
afzonderlijk worden behandeld; hoe meer primitieve symbolen, des te
langer onze bewijzen.
\end{explain}
\end{tagblock}

% Alternate symbols

\begin{intro}
Mogelijk kent u andere termen en symbolen dan de bovenstaande. In
logicaboeken en in het onderwijs gebruikt men vaak $\sim$, $\neg$ of~!\ voor
``negatie'', en $\wedge$, $\cdot$ of $\&$ voor ``conjunctie''.
Veelgebruikte symbolen voor de ``conditionele verbinding'' of
``implicatie'' zijn $\rightarrow$, $\Rightarrow$ en $\supset$.
\iftag{prvIff,defIff}{Symbolen voor ``biconditionele verbinding'', ``bi-implicatie''
  of ``(materiële) equivalentie'' zijn $\leftrightarrow$,
  $\Leftrightarrow$ en $\equiv$.}{}
\iftag{prvFalse,defFalse}{Het symbool $\lfalse$ heet ook wel
  ``onwaarheid'', ``falsum'', ``absurdum'' of ``bottom''.}{}
\iftag{prvTrue,defTrue}{Het symbool $\ltrue$ heet ook wel
  ``waarheid'', ``verum'' of ``top''.}{}

Gewoonlijk gebruikt men kleine letters aan het begin van het Latijnse
alfabet (bijvoorbeeld $a$, $b$, $c$) voor !!{constant}s, die soms namen
heten, en kleine letters aan het einde (bijvoorbeeld $x$, $y$, $z$)
voor !!{variable}s. Kwantoren worden met !!{variable}s gecombineerd,
bijvoorbeeld met $x$. Notatievarianten voor de universele kwantor zijn
$\forall x$, $(\forall x)$, $(x)$, $\Pi x$, $\bigwedge_x$; voor de
existentiële kwantor zijn dat $\exists x$, $(\exists x)$, $(Ex)$,
$\Sigma x$, $\bigvee_x$.
\end{intro}

\begin{explain}
We kunnen alle propositionele operatoren en beide kwantoren als
primitieve symbolen van de taal behandelen. We kunnen ook een kleinere
voorraad primitieve symbolen kiezen en de overige !!{operator}s als
gedefinieerd behandelen. Functioneel volledige verzamelingen van
Booleaanse operatoren zijn onder meer $\{ \lnot, \lor \}$,
$\{ \lnot, \land \}$ en $\{ \lnot, \lif\}$. Elk daarvan kan met een
van beide kwantoren worden gecombineerd tot een expressief volledige
eerste-ordetaal.

Mogelijk kent u nog twee andere !!{operator}s: de Sheffer-streep~$|$,
genoemd naar Henry Sheffer, en de pijl van Peirce~$\downarrow$, ook
bekend als de dolk van Quine. Met hun gebruikelijke lezingen ``nand''
respectievelijk ``nor'' zijn deze operatoren elk afzonderlijk
functioneel volledig.
\end{explain}
\end{document}
